% ---------
% --------
%!TEX TS-program = pdflatex
%!TEX encoding = UTF-8 Unicode
%!TEX spellcheck = English (Aspell)

\documentclass[11pt]{article}
\usepackage{jeffe,handout,graphicx}
\usepackage{fourier-orns}
\usepackage{mathtools}
\usepackage[normalem]{ulem} % either use this (simple) or
\usepackage{tikz}


\def\Sym#1{\texttt{\color{BrickRed}#1}}
\allowdisplaybreaks

\begin{document}

\begin{center}
\Large\textbf{CSCI 432/532, Spring 2026}%
\\
\LARGE\textbf{Homework 7}%
\\[0.5ex]
\large Due Monday, March 9 at 8:30am
\end{center}

\bigskip
\hrule
\bigskip

\subsection*{Submission Requirements}
\begin{itemize}
    \item Type or clearly hand-write your solutions into a PDF format so that they are legible and professional. Submit your PDF on Canvas or Gradescope. 
	\item Put each \textbf{sub}problem on a separate page and select the corresponding problem via the Gradescope interface.
    \item Do not submit your first draft. Type or clearly re-write your solutions for your final submission. If your submission is not legible, we will ask you to resubmit.
\end{itemize}

\subsection*{Collaboration and Resources}

The \emph{best} resources for completing the homework are (in no particular order):
\begin{itemize}
	\item Your own notes from lectures and problem sessions, and/or lecture recordings.
	\item The lecture notes linked from the course website for the lecture day.
	\item The questions channel on Discord.
	\item Office hours.
	\item Discussions with classmates.
\end{itemize}
You may also access almost any resource in preparing your solution to the
homework. The exception is that you may not seek answers to the problems on the
homework directly, such as by pasting them into a GenAI tool, searching for
solutions on a search engine, looking for them on Chegg or similar websites, or
asking another person for the solution.

Additionally, you \EMPH{must}
\begin{itemize}
    \item Write your solutions in your own words---this means that you should never be \emph{copying}
    anything of substance from any source, and
    \item credit every resource you use, including GenAI tools, by providing links or full chat transcripts.
    If you use the provided LaTeX template, you can use the \texttt{Sources}
    environment for this. Ask if you need help!
\end{itemize}


\begin{boxquote}{0.9}
Remember, if you choose to use GenAI tools, not only must you provide a full
transcript of your conversation (or a link to one), you must also use the
following prompt (or some variation of it), and provide a full transcript of
the conversation.
\end{boxquote}  

Prompt:
\begin{boxquote}{0.9}
You are acting as a teaching assistant for an advanced undergraduate/graduate
algorithms and computer science theory course.

Your role is not to provide full solutions or final answers.

Instead, you should act like we are having a conversation. Ask me a single
question and let me respond. Avoid asking me many questions at once or giving
me a lot of information at one time. Guide me using hints, leading questions,
partial insights, and sanity checks.  Help me debug my own proofs, algorithms,
or reductions. Do not give a complete solution or full proofs. Be patient with
me. Don't give me details so that I can move on to the next part of a problem.
Make sure that I understood before moving on.

The course uses Jeff Erickson's Models of Computation lecture notes, so you should
guide me to answer problems using the definitions, notation, and style from those notes.

Assume I am a serious student trying to learn, not to shortcut the assignment.
I am attaching the assignment, so you should refuse to provide a complete solution
to the problems listed. Of course, you can walk me through the "Solved
Problems" listed at the end if I ask.
\end{boxquote}  

\newpage
%----------------------------------------------------------------------

\headers{CSCI 432/532}{Homework 7 (due March 9)}{Spring 2026}

\begin{enumerate}
%\parindent 1.5em \itemsep 4ex plus 0.5fil

%----------------------------------------------------------------------

\item 
Recall the language $\mathsc{NeverReject} = \{ \langle M \rangle :
\textsc{Reject}(M) = \emptyset \}$ from last homework. 
 Prove that $\mathsc{NeverReject}$ is undecidable via \EMPH{reduction}.

\item
Let $\seq{M}$ denote the string encoding of a Turing machine $M$. Recall that if
$w$ is a string, then $ww$ is a new string that repeats $w$ twice.  Prove that
the following language is undecidable. 
\[
	\mathsc{SelfRepeatAccept} 
	:=
	\Setbar{\seq{M}}{\text{$M$ accepts the string $\seq{M} \seq{M}$}}
\]

\begin{Rubric}{Undecidability proof rubric}~

\begin{itemize}

\item \textbf{Diagonalization}
\begin{itemize}
\item[+ 4] for a correct wrapper Turing machine
\item[+ 6] for self-contradiction proof (= 3 for $\Longleftarrow$ + 3 for $\Longrightarrow$)
\end{itemize}

\item \textbf{Reduction}
\begin{itemize}
\item[+ 4] for a correct reduction
\item[+ 3] for ''if" proof
\item[+ 3] for ''only if" proof
\end{itemize}

\end{itemize}
\end{Rubric}

\item
Recall that a \EMPH{vertex cover} of a graph is a set of a vertices that touches every edge of that graph.
Suppose you are given a magic black box that somehow answers the following decision problem \emph{in polynomial time}:
\begin{itemize}
    \item \textsc{Input}: an undirected graph $G$ and a positive integer $k$.
    \item \textsc{Output}: \textsc{True} if $G$ has a vertex cover of size $k$ and \textsc{False} otherwise.
\end{itemize}

\begin{enumerate}[(a)]
    \item Using this black box as a subroutine, describe an algorithm that
    solves the following \EMPH{optimization problem} \emph{in polynomial time}:
\begin{itemize}
    \item \textsc{Input}: an undirected graph $G$.
    \item \textsc{Output}: the size of the smallest vertex cover of $G$.
\end{itemize}

\item Using this black box as a subroutine, describe an algorithm that solves
the following \EMPH{search problem} \emph{in polynomial time}:
\begin{itemize}
    \item \textsc{Input}: an undirected graph $G$.
    \item \textsc{Output}: a vertex cover in $G$ of smallest size.
\end{itemize}
\end{enumerate}

\emph{Unlike in the problem session solutions}, you do not need to argue the correctness of your algorithms. But you still do need to argue that they run in polynomial time.

\item
Formally, a \EMPH{proper coloring} of a graph $G=(V,E)$ is a function $c\colon V \to \set{1,2,\dots,k}$, for some integer $k$, such that $c(u)\ne c(v)$ for all $uv\in E$.  Less formally, a valid coloring assigns each vertex of $G$ a color, such that every edge in $G$ has endpoints with different colors.  The \EMPH{chromatic number} of a graph is the minimum number of colors in a proper coloring of $G$.

Suppose you are given a magic black box that somehow answers the following
decision problem \emph{in polynomial time}:
\begin{itemize}
\item \mathsc{Input:} An undirected graph $G$ and an integer $k$.
\item \mathsc{Output:} \mathsc{True} if $G$ has a proper coloring with $k$ colors, and \mathsc{False} otherwise.
\end{itemize}
Using this black box as a subroutine, describe an algorithm that solves the
following \EMPH{search problem} \emph{in polynomial time}:
\begin{itemize}
\item \mathsc{Input:} An undirected graph $G$.
\item \mathsc{Output:} A valid coloring of $G$ using the minimum possible number
of colors.
\end{itemize}
\Hint{You can use the magic box more than once.  The input to the magic box is a graph and \textbf{only} a graph, meaning \textbf{only} vertices and edges.}

\emph{Unlike in the problem session solutions}, you do not need to argue the
correctness of your algorithms. But you still do need to argue that they run in
polynomial time.

\end{enumerate}

\begin{Rubric}{Basic reduction (algorithms)}
10 points =
\begin{itemize}
\item[+ 4] for calling the magic black box as a subroutine in the algorithm. No points here if it is not called on the correct input data type.
\item[+ 2] for a correct algorithm.
\item[+ 2] for an algorithm that runs in polynomial time (assuming the black box runs in polynomial time).
\item[+ 2] for a valid argument that the algorithm runs in polynomial time.
\end{itemize}
\end{Rubric}



% ===========================================================================
\newpage

\subsection*{Solved problems}

\begin{enumerate}\parindent 1.5em

\item
Consider the language $\mathsc{SometimesHalt} = \set{ \seq{M} \mid \text{$M$~halts
on at least one input string}}$.  Note that $\seq{M}\in \mathsc{SometimesHalt}$ does not imply that $M$ \emph{accepts} any strings; it is enough that $M$ \emph{halts} on (and possibly rejects) some string.

Prove that \mathsc{SometimesHalt} is undecidable.

\begin{Solution}
We can reduce the standard halting problem to \mathsc{Sometimes\-Halt} as follows:
\begin{algo}
	\textul{$\mathsc{DecideHalt}(\seq{M, w})$:}\+
\\	Encode the following Turing machine $M'$:\+
\\	\color{OliveGreen}
	\qquad \begin{algorithm}
		\textul{$M'(x)$:}\+
	\\		(ignore $x$)
	\\		run $M$ on input $w$
	\end{algorithm}
\\[1ex]
	return $\mathsc{DecideSometimesHalt}(\seq{M'})$\-
\end{algo}
We prove this reduction correct as follows:

\begin{itemize}
\item[$\Longrightarrow$]
Suppose $M$ halts on input $w$.\\
Then $M'$ halts on \emph{every} input string $x$.\\
So \mathsc{DecideSometimesHalt} must accept the encoding $\seq{M’}$.\\
We conclude that \mathsc{DecideHalt} correctly accepts the encoding $\seq{M, w}$.

\item[$\Longleftarrow$]
Suppose $M$ does not halt on input $w$.\\
Then $M'$ diverges on \emph{every} input string $x$.\\
So \mathsc{DecideSometimesHalt} must reject the encoding $\seq{M’}$.\\
We conclude that \mathsc{DecideHalt} correctly rejects the encoding $\seq{M, w}$.
\end{itemize} 
\end{Solution}

\end{enumerate}

\end{document}

